\section{VulnRAG: System Architecture}
\label{sec:design}

VulnRAG is designed as a stateful, modular framework that decouples vulnerability reasoning from raw data retrieval. Unlike traditional monolithic RAG pipelines, our architecture (Fig.~\ref{fig:arch}) utilizes a Master-Worker hierarchy to coordinate specialized agents across heterogeneous LLM backends.

\begin{figure}[htbp]
    \centering
    \resizebox{\columnwidth}{!}{\begin{tikzpicture}[node distance=1.5cm, every node/.style={font=\footnotesize}]
    % Styles
    \tikzstyle{block} = [rectangle, rounded corners, minimum width=2.2cm, minimum height=0.8cm, text centered, draw=black, fill=blue!10, drop shadow]
    \tikzstyle{database} = [cylinder, shape border rotate=90, draw, minimum width=1.5cm, minimum height=1.8cm, text centered, fill=gray!20, drop shadow]
    \tikzstyle{agent} = [rectangle, rounded corners, minimum width=2.0cm, minimum height=0.7cm, text centered, draw=black, fill=green!10, drop shadow]
    \tikzstyle{arrow} = [thick,->,>=stealth]

    % Nodes
    \node (nvd) [block, fill=orange!20] {NVD API 2.0};
    \node (sync) [block, right=of nvd] {NVD Sync Service};
    \node (qdrant) [database, right=of sync] {Qdrant Vector Store};
    
    \node (query) [above=of sync, xshift=-1cm] {User Security Query};
    \node (reform) [agent, above=of sync, xshift=1cm] {Query Reformulator};
    
    \node (search) [agent, right=of reform] {Vector Search};
    \node (rerank) [agent, right=of search] {MFSR Reranker};
    \node (eval) [agent, above=of search, fill=red!10] {Evaluator Agent};

    % Connections
    \draw [arrow] (nvd) -- (sync);
    \draw [arrow] (sync) -- (qdrant);
    
    \draw [arrow] (query) -- (reform);
    \draw [arrow] (reform) -- (search);
    \draw [arrow] (qdrant) -- (search);
    \draw [arrow] (search) -- (rerank);
    \draw [arrow] (rerank) -- (eval);
    
    % Feedback Loop
    \draw [arrow, dashed] (eval) -| node[above, pos=0.25] {Self-Correction Feedback} (reform);
    
    % Output
    \node (out) [block, right=of rerank, fill=cyan!10] {Intelligence Synthesis};
    \draw [arrow] (rerank) -- (out);

    % Legend or Labels
    \node [below=0.3cm of qdrant, font=\tiny] {BGE-Large Embeddings};
    \node [below=0.3cm of rerank, font=\tiny] {Cross-Encoder + Domain Metrics};

\end{tikzpicture}
}
    \caption{VulnRAG System Architecture and Information Flow.}
    \label{fig:arch}
\end{figure}

\subsection{Knowledge Synchronization and Storage}
The foundation of VulnRAG is a continuously updated knowledge base of vulnerability disclosures. 
\begin{itemize}
    \item \textbf{Continuous Sync}: The \texttt{nvd\_sync\_service.py} module manages real-time synchronization with the NIST NVD API 2.0. It parses raw JSON advisories into a structured format, normalizing software versioning (CPE) and security metrics (CVSS).
    \item \textbf{Vector Store}: We utilize \textbf{Qdrant} for high-performance vector storage. CVE descriptions and metadata are embedded using the \textbf{BAAI/bge-large-en-v1.5} model via the \texttt{SentenceTransformers} framework.
    \item \textbf{Metadata Filtering}: To improve retrieval precision, we implement hybrid search patterns in \texttt{cve\_vector\_store.py} that combine semantic embeddings with payload-based metadata filters (e.g., CVSS severity, vendor tags, and exploit availability).
\end{itemize}

\subsection{Master-Worker Orchestration}
The system's reasoning logic is managed by the \texttt{VulnIntelAgent}, which orchestrates the following specialized sub-modules:
\begin{enumerate}
    \item \textbf{Query Reformulator}: Enhances raw user queries with domain-specific terminology.
    \item \textbf{Smart Searcher}: Coordinates the retrieval and reranking loop.
    \item \textbf{Graph Explorer}: Discovers structural relationships between vulnerabilities.
    \item \textbf{Evaluator}: Assesses result quality and triggers self-correction loops.
\end{enumerate}

By modularizing these components, VulnRAG can adapt its reasoning strategy based on the complexity of the security task at hand.
